\section{The exact parameter certificate}
\label{app:certificate}

The following Python program uses only integer and rational arithmetic.
It implements the prescription in Section~\ref{sec:parameters}, checks
every parameter interval and the endpoint derivative inequalities, and
verifies the stopping indices and gaps in
Table~\ref{tab:parameter-certificate}.

\begingroup\footnotesize
\begin{verbatim}
from fractions import Fraction as F
from math import comb

K, G = 3000, 100000

def polynomials(d):
    k = d - 2
    for i in range(1, 500):
        t = F(i, 50)
        values = [sum(comb(d-r, j-s) * comb(k, j-1)
                      * k**(j-1) * t**(j-s)
                      for j in range(1, d))
                  for r, s in [(0, 0), (1, 0), (1, 1)]]
        yield (t, *values)

def upper(x, e):
    assert 0 <= x < 1
    l, r = 0, G
    while r > l + 1:
        middle = (l + r) // 2
        if F(middle, G)**e > x:
            r = middle
        else:
            l = middle
    return F(r, G)

def check_derivatives():
    for k in range(2, 7):
        for t in [F(3, 20), F(5)]:
            for a, b in [(1, 0), (0, 1), (1, 1)]:
                for h in range(2, k+3):
                    N = t * (1 + 1/t
                        - F(h-1, h)*(k+1)/(k+t)
                        - F(1, h)*(a+b)/(a+b*t))
                    assert N >= 0
                    assert N**h < (k+t)**(h-1)*(a+b*t)

def check_degree(d):
    k, choices, least = d-2, list(polynomials(d)), F(1)
    for z in range(K, 2*K):
        etaL, etaH = F(z, K), F(z+1, K)
        if etaL*(1-k*etaL/d) < F(2495, 10000):
            return z, least
        bL, bH = d-k*etaH, d-k*etaL
        t, H, Ha, Hb = min(choices,
            key=lambda row: (etaH+bH*row[0])**d/row[1])
        w = upper(((d+1)*(etaH+bH*t)/d**2)**d/H, d-1)
        assert (w**(d-2)*min(Ha, Hb)
                > ((etaH+bH*t)/(d-1))**(d-1))
        m = (d-1)*w
        u, v = (etaH-m)/d, (bH-m)/d
        assert 0 < m < 1 and min(etaL, bL) > m
        assert min(u, v) > 0 and F(3, 20) <= v/u <= 5

        def delta(h, D):
            return upper((u*v)**h/((k*u+v)**(h-1)*D), h)

        P = [delta(d, u+v)]
        if d == 4 and z < 4000:
            assert m+2*u < F(999, 1000)
            P += [delta(3, v),
                  max(delta(2, min(3*u, 2*v, u+v)),
                      upper(u*u/3, 2))]
        else:
            P += [delta(h, min(u, v))
                  for h in range(d-1, 1, -1)]
        g = (d-1)*(k-1) + (d == 4)
        P.append(max(u/2, min(u, v/g)))
        gap = 1-m-sum(max(P[:h]) for h in range(1, d+1))
        assert gap > F(1, 1000)
        least = min(least, gap)
    raise AssertionError("cutoff not reached")

check_derivatives()
assert [check_degree(d) for d in range(4, 9)] == [
    (5124, F(191, 100000)), (4084, F(293, 50000)),
    (3552, F(321, 2500)), (3226, F(23177, 100000)),
    (3003, F(31683, 100000))]
\end{verbatim}
\endgroup

\subsection{The enlarged first quartic cell}
\label{app:quartic-extension}

The following additional exact checks verify the parameters used in
Proposition~\ref{prop:cyclotomic-central}. They use the functions
\texttt{upper} and \texttt{F} from the program above. The fixed
bidegree ratio is $t_*=1/2$, for which
$(H,H_a,H_b)=(10,5,10)$. In particular, no continuity argument or
floating-point computation is needed for the small extension below
$\eta=1$.

\begingroup\footnotesize
\begin{verbatim}
epsilon = F(1, 10**6)
etaL, etaH = 1-epsilon, F(3001, 3000)
bL, bH = 4-2*etaH, 2+2*epsilon
t, w = F(1, 2), F(24809, 100000)
m = 3*w
u, v = (etaH-m)/4, (bH-m)/4
assert 0 < epsilon < F(1, 2)
assert 0 < m < 1 and m < min(etaL, bL)
assert w**3*10 > (F(5, 16)*(etaH+bH*t))**4
assert w**2*5 > ((etaH+bH*t)/3)**3
assert min(u, v) > 0 and F(3, 20) <= v/u <= 5
def quartic_delta(h, D):
    return upper((u*v)**h/((2*u+v)**(h-1)*D), h)
P = [quartic_delta(4, u+v), quartic_delta(3, v),
     max(quartic_delta(2, min(3*u, 2*v, u+v)),
         upper(u*u/3, 2)), max(u/2, min(u, v/4))]
assert P == [F(4729, 100000), F(5097, 100000),
             F(6899, 100000), F(76819, 1200000)]
gap = 1-m-sum(max(P[:i]) for i in range(1, 5))
tau = F(1, 50000)
assert gap-4*tau == F(1941, 100000)
assert 1-m-2*u-2*tau == F(15319, 120000)
\end{verbatim}
\endgroup
